% =============================================================
%  Resumo do Projeto da Dissertação (1,00 ponto)
% =============================================================
\section*{Resumo do Projeto da Dissertação}

A expansão de sistemas de monitoramento urbano tornou o vídeo uma fonte cada vez mais relevante de prova material em investigações criminais e cíveis. Diferentemente de câmeras dedicadas ao controle de velocidade, câmeras de monitoramento de cidades não foram projetadas para fins periciais: têm posicionamento, resolução, ângulo e taxa de quadros variáveis, o que torna sua análise tecnicamente mais complexa. O trabalho do perito diante desse material envolve localizar trechos relevantes em longas gravações, capturar e organizar frames, interpretar o conteúdo visual e produzir uma demonstração tecnicamente fundamentada para o laudo --- processo hoje mais custoso, em tempo e esforço, do que qualquer cálculo isolado.

Paralelamente, algoritmos de processamento de vídeo (detecção e rastreamento de objetos, reconhecimento de eventos, detecção de anomalias) e modelos multimodais de linguagem já demonstram capacidade de descrever e raciocinar sobre conteúdo visual em linguagem natural, abrindo a possibilidade de um sistema de apoio que combine algoritmos especializados de visão computacional com um modelo de linguagem.

Este projeto propõe, portanto, realizar uma análise comparativa de diferentes algoritmos que possam complementar um modelo de linguagem no apoio ao trabalho do perito criminal na análise de vídeos de câmeras de monitoramento urbano --- da captura dos frames relevantes à geração de uma análise preliminar organizada. Metodologicamente, prevê-se revisão de literatura, definição de critérios de comparação, testes com um conjunto de vídeos representativos e avaliação da viabilidade prática de uso na rotina pericial. Espera-se, como resultado, uma caracterização comparativa dos algoritmos testados e um protocolo/arranjo de referência aplicável a unidades de perícia criminal.
